\documentclass{standalone}
\usepackage{tikz}
\usepackage{luacode}

\def\escapetime#1#2{\directlua{tex.print(escape_time(#1,#2))}}
\def\mandcolor#1#2{\directlua{mand_color(#1,#2)}}

\begin{document}
\begin{luacode}
local NUMITER=100

point_iter = function (a,b) 
   local n = 1
   local x = a
   local y = b
   return function ()
      if n > NUMITER then return nil end
      n = n+1
      x,y = x*x - y*y + a, 2*x*y + b
      return n, x*x + y*y
   end
end

escape_time = function (a,b)
   local esc_time
   for n, d in point_iter(a,b) do
      esc_time = n
      if d > 4 then break end
   end
   return esc_time
end

mand_color = function (a,b)
   local etime = escape_time(a,b)
   if etime == NUMITER+1 then 
      tex.print("black") 
   else 
      local percent = 100 * etime / NUMITER
      tex.print("red!"..percent)
   end
end
\end{luacode}

\begin{tikzpicture}[scale=2]
   \foreach \x in {-2,-1.99,...,.5}{
      \foreach \y in {-1.3,-1.29,...,1.3}{
      \draw[\mandcolor{\x}{\y},fill] (\x,\y) +(-.005,-.005) rectangle +(.005,.005);}
   }
 
 %dibujo los rayos entorno al 5 ciclo
 \draw   (-0.5,0.5) --(-0.65,0.55);  \draw  (-0.65,0.55)--(-0.8,0.85);
 \draw     (-0.47,0.5) --(-0.45,0.65);  \draw  (-0.45,0.65)--(-0.6,0.95);
 
 %dibujo flechas y etiqueta
 \draw[<->]  (-0.8,0.85)--(-0.6,0.95);
 \node at (-0.75,1.05) {\scriptsize$\frac{1}{31}$};
\end{tikzpicture}
\end{document}